\chapter{Patologia si recuperarea post-AVC}

\section{Patologia accidentului vascular cerebral}
\hspace{0.75cm}Accidentul vascular cerebral reprezintă un deficit neurologic focal cu instalare bruscă, cauzat de o perturbare severă 
a circulației sanguine la nivelul encefalului. Din punct de vedere fiziopatologic, AVC-ul se clasifică în două categorii 
majore: ischemic (care reprezintă aproximativ 85\% din cazuri, cauzat de ocluzia unui vas de sânge prin tromboză sau embolism)
și hemoragic (cauzat de ruptura vasculară) \cite{stroke_patho_2020}. Indiferent de etiologie, privarea țesutului cerebral de
oxigen și glucoză declanșează o cascadă ischemică distructivă. În centrul leziunii neuronii 
necrozeaza în câteva minute. Supraviețuirea pacientului este adesea acompaniată
de deficite motorii severe pe partea opusă emisferei cerebrale afectate (hemipareză sau hemiplegie), tulburări de echilibru 
și spasticitate musculară, care compromit drastic independența funcțională a acestuia.

\section{Mecanismul biologic al recuperării: neuroplasticitatea}

\hspace{0.75cm}Spre deosebire de dogma medicală clasică, care susținea că sistemul nervos central adult este o structură rigidă, cercetările
moderne au demonstrat că recuperarea funcțiilor motorii pierdute este posibilă datorită neuroplasticității. Aceasta este definită ca 
abilitatea intrinsecă a creierului de a se remodela structural și funcțional ca răspuns la experiență, învățare sau leziuni \cite{kleim_jones_2008}.
După un AVC, creierul inițiază un proces de reorganizare corticală: neuronii din zonele adiacente leziunii, sau chiar din emisfera 
contralaterală (sănătoasă), pot prelua parțial funcțiile rețelelor neuronale distruse. Acest fenomen implică formarea de noi sinapse 
(sinaptogeneză), modificarea eficienței sinapselor existente și ramificarea dendritică.

\hspace{0.75cm}La nivel național și internațional, contribuții fundamentale la aprofundarea acestor mecanisme
celulare și a capacității creierului de autoreparare 
au fost aduse de academicianul și neurochirurgul român Prof. Dr. Leon Dănăilă \cite{danaila_2023}.
Desi structura cerebrală dispune de acest instrumentar celular uimitor de regenerare, neuroplasticitatea rămâne un proces
„dependent de experiență”; ea nu se produce spontan
la un nivel suficient de înalt pentru a reda independența funcțională, ci trebuie stimulată activ
și susținut prin reînvățare motorie \cite{cramer_2011}.

\section{Importanța mișcărilor repetitive}
\hspace{0.75cm}Pentru a transforma potențialul neuroplasticității în recuperare funcțională reală, reabilitarea motorie trebuie să respecte 
anumite principii biomecanice și neurologice stricte. Unul dintre cele mai importante este principiul repetiției. Studiile de neurofiziologie 
arată că sunt necesare mii de repetări ale unei mișcări specifice pentru a consolida noile căi neuronale și a transforma o mișcare stângace, 
conștientă, într-un automatism motor \cite{lang_2009}. Terapia indusă prin constrângere (CIMT - Constraint-Induced Movement Therapy) 
și antrenamentul specific sarcinii (task-specific training) demonstrează că intensitatea și frecvența ridicată a exercițiilor dictează 
gradul de recuperare. În lipsa unor mișcări repetitive, ghidate corect, pacientul riscă să dezvolte mișcări compensatorii vicioase 
(folosind doar membrele sănătoase), 
fenomen cunoscut sub numele de „neutilizare învățată” (learned non-use), care plafonează recuperarea pe termen lung \cite{taub_2014}.
